This section briefly characterizes the \cmap\ data layer at the level needed to contextualize the agentic system; a fuller account of the underlying portal is given in \citet{Ashkezari2021SimonsCMAP}.
\cmap\ harmonizes heterogeneous datasets into a common structure and indexes entries by location and time, enabling efficient retrieval and integration. The system exposes these capabilities via interactive and programmatic interfaces.

\paragraph{Curation and harmonization as a substrate for agents.}
CMAP is assembled through domain-scientist curation: ingest and documentation practices standardize metadata, references, and variable semantics across heterogeneous sources.
For an agentic interface, this curated substrate has two concrete consequences.
First, retrieval can anchor decisions in vetted descriptors---units, spatial and temporal coverage, instrument or model provenance---so that the agent's choice of a dataset or variable is grounded in catalog facts rather than inferred from the model's latent knowledge.
Second, tools can operate on consistent identifiers (table and variable short names) without per-dataset special cases, which is what makes a single, generic execution layer practical across hundreds of products.
We therefore treat the curated layer as a scientific enabler of reliable discovery and integration, not merely as an implementation convenience.
